\documentclass[10pt,a4paper]{amsart}
\usepackage{amsmath,amssymb,amsthm,graphicx}
% Retained minimal preamble for the Stacks Project chapter more-morphisms.tex (More on Morphisms).
% Verbatim excerpt of the chapter preamble: Xy-pic setup (native preamble lines 14-18) and the native
% theorem declarations with the mathematical macros (native preamble lines 154-217). Copied mechanically;
% nothing added, removed or redefined. The wrapper class and the amsmath/amssymb/amsthm/graphicx base
% packages are supplied by the assembler, so the class/externaldocument part of the native preamble is omitted.
% For commutative diagrams we use Xy-pic
\usepackage[all]{xy}

% We use 2cell for 2-commutative diagrams.
\xyoption{2cell}
% Theorem environments.
%
\theoremstyle{plain}
\newtheorem{theorem}[subsection]{Theorem}
\newtheorem{proposition}[subsection]{Proposition}
\newtheorem{lemma}[subsection]{Lemma}

\theoremstyle{definition}
\newtheorem{definition}[subsection]{Definition}
\newtheorem{example}[subsection]{Example}
\newtheorem{exercise}[subsection]{Exercise}
\newtheorem{situation}[subsection]{Situation}

\theoremstyle{remark}
\newtheorem{remark}[subsection]{Remark}
\newtheorem{remarks}[subsection]{Remarks}

\numberwithin{equation}{subsection}

% Macros
%
\def\lim{\mathop{\mathrm{lim}}\nolimits}
\def\colim{\mathop{\mathrm{colim}}\nolimits}
\def\Spec{\mathop{\mathrm{Spec}}}
\def\Hom{\mathop{\mathrm{Hom}}\nolimits}
\def\Ext{\mathop{\mathrm{Ext}}\nolimits}
\def\SheafHom{\mathop{\mathcal{H}\!\mathit{om}}\nolimits}
\def\SheafExt{\mathop{\mathcal{E}\!\mathit{xt}}\nolimits}
\def\Sch{\mathit{Sch}}
\def\Mor{\mathop{\mathrm{Mor}}\nolimits}
\def\Ob{\mathop{\mathrm{Ob}}\nolimits}
\def\Sh{\mathop{\mathit{Sh}}\nolimits}
\def\NL{\mathop{N\!L}\nolimits}
\def\CH{\mathop{\mathrm{CH}}\nolimits}
\def\proetale{{pro\text{-}\acute{e}tale}}
\def\etale{{\acute{e}tale}}
\def\QCoh{\mathit{QCoh}}
\def\Ker{\mathop{\mathrm{Ker}}}
\def\Im{\mathop{\mathrm{Im}}}
\def\Coker{\mathop{\mathrm{Coker}}}
\def\Coim{\mathop{\mathrm{Coim}}}

% Boxtimes
%
\DeclareMathSymbol{\boxtimes}{\mathbin}{AMSa}{"02}

%
% Macros for moduli stacks/spaces
%
\def\QCohstack{\mathcal{QC}\!\mathit{oh}}
\def\Cohstack{\mathcal{C}\!\mathit{oh}}
\def\Spacesstack{\mathcal{S}\!\mathit{paces}}
\def\Quotfunctor{\mathrm{Quot}}
\def\Hilbfunctor{\mathrm{Hilb}}
\def\Curvesstack{\mathcal{C}\!\mathit{urves}}
\def\Polarizedstack{\mathcal{P}\!\mathit{olarized}}
\def\Complexesstack{\mathcal{C}\!\mathit{omplexes}}
% \Pic is the operator that assigns to X its picard group, usage \Pic(X)
% \Picardstack_{X/B} denotes the Picard stack of X over B
% \Picardfunctor_{X/B} denotes the Picard functor of X over B
\def\Pic{\mathop{\mathrm{Pic}}\nolimits}
\def\Picardstack{\mathcal{P}\!\mathit{ic}}
\def\Picardfunctor{\mathrm{Pic}}
\def\Deformationcategory{\mathcal{D}\!\mathit{ef}}
\newenvironment{reference}{}{}
\newenvironment{slogan}{}{}
\newenvironment{history}{}{}
\usepackage{cleveref}
\crefname{originallemma}{original lemma}{original lemmas}
\Crefname{originallemma}{Original lemma}{Original lemmas}
\makeatletter
\newlabel{lemma-relative-pseudo-coherence-characterize}{{09UJ}{source}{}{sourcepointer-1}{}}
\newlabel{lemma-relative-pseudo-coherence-characterize@cref}{{[originallemma][0][]09UJ}{[1][0][]source}{}{}{}}
\makeatother
\crefname{originallemma}{original lemma}{original lemmas}
\Crefname{originallemma}{Original lemma}{Original lemmas}
\makeatletter
\newlabel{morphisms-lemma-homeomorphism-affine}{{04DE}{source}{}{sourcepointer-2}{}}
\newlabel{morphisms-lemma-homeomorphism-affine@cref}{{[originallemma][0][]04DE}{[1][0][]source}{}{}{}}
\makeatother
\crefname{originallemma}{original lemma}{original lemmas}
\Crefname{originallemma}{Original lemma}{Original lemmas}
\makeatletter
\newlabel{perfect-lemma-closed-push-pseudo-coherent}{{09VA}{source}{}{sourcepointer-3}{}}
\newlabel{perfect-lemma-closed-push-pseudo-coherent@cref}{{[originallemma][0][]09VA}{[1][0][]source}{}{}{}}
\makeatother
\title{A universally closed-on-points map over the base detects relative pseudo-coherence through its pushforward}
\author{The Stacks Project Authors}
\date{stacks-project:more-morphisms.tex:local-snapshot-20261011; GNU FDL 1.2 or later}
\begin{document}
\maketitle
\noindent\textit{Editorial note.} Statement span is lines 16948-16955 of sources/stacks-more-morphisms/pinned-original-source.tex; proof span is lines 16957-16992, the complete author proof (34 body lines).\par\medskip
\noindent\textit{Editorial note.} Navigation stubs: 3 cited-but-not-included labels are supplied as non-load-bearing pointers with their pinned tags where the snapshot tag index has them.\par\medskip
\noindent\textit{Editorial note.} extra\_packages is empty on purpose: the retained preamble already loads xy with options and a second declaration clashes in the assembler wrapper.\par\medskip
\noindent\textit{Editorial note.} Text is retained verbatim, including the author's own typos and placeholders.\par\medskip
\noindent\textit{Editorial note.} The author writes rad for the Jacobson radical in this proof; retained verbatim.\par\medskip
\noindent\textit{Allowed standard prerequisites (editorial).}\begin{itemize}
\item Explicit prerequisite boundary. Affineness of point-closed maps locally of finite type (Morphisms, lemma-homeomorphism-affine), the affine-space test for relative pseudo-coherence (same chapter), and invariance under closed pushforwards (Derived Categories of Schemes, lemma-closed-push-pseudo-coherent).
\end{itemize}\medskip
\ifdefined\hypertarget\hypertarget{sourcepointer-1}{}\fi
\noindent\textit{Source reference.} Original lemma 09UJ: sources/stacks-more-morphisms/pinned-original-source.tex, source lines 16837--16847. Cited by the retained statement or proof but not part of the collected spans; supplied only as a navigation pointer. The argument does not depend on its content.\par\medskip
\ifdefined\hypertarget\hypertarget{sourcepointer-2}{}\fi
\noindent\textit{Source reference.} Original lemma 04DE: sources/stacks-more-morphisms/pinned-original-source.tex, source lines 16992--16992. Cited by the retained statement or proof but not part of the collected spans; supplied only as a navigation pointer. The argument does not depend on its content.\par\medskip
\ifdefined\hypertarget\hypertarget{sourcepointer-3}{}\fi
\noindent\textit{Source reference.} Original lemma 09VA: sources/stacks-more-morphisms/pinned-original-source.tex, source lines 16992--16992. Cited by the retained statement or proof but not part of the collected spans; supplied only as a navigation pointer. The argument does not depend on its content.\par\medskip
\begin{lemma}
\label{lemma-closed-morphism-relative-pseudo-coherence}
Let $i : X \to Y$ morphism of schemes locally of finite type over a
base scheme $S$. Assume that $i$ induces a homeomorphism of $X$ with a closed
subset of $Y$. Let $E$ be an object of $D(\mathcal{O}_X)$.
Then $E$ is $m$-pseudo-coherent relative to $S$ if and only if
$Ri_*E$ is $m$-pseudo-coherent relative to $S$.
\end{lemma}

\begin{proof}
By Morphisms, Lemma \ref{morphisms-lemma-homeomorphism-affine}
the morphism $i$ is affine. Thus we may assume $S$, $Y$, and $X$ are affine.
Say $S = \Spec(R)$, $Y = \Spec(A)$, and $X = \Spec(B)$.
The condition means that $A/\text{rad}(A) \to B/\text{rad}(B)$ is
surjective; here $\text{rad}(A)$ and $\text{rad}(B)$ denote
the Jacobson radical of $A$ and $B$.
As $B$ is of finite type over $A$, we can find
$b_1, \ldots, b_m \in \text{rad}(B)$ which generate $B$ as an
$A$-algebra. Say $b_j^N = 0$ for all $j$. Consider the diagram of
rings
$$
\xymatrix{
B & R[x_i, y_j]/(y_j^N) \ar[l] & R[x_i, y_j] \ar[l] \\
A \ar[u] & R[x_i] \ar[l] \ar[u] \ar[ru]
}
$$
which translates into a diagram
$$
\xymatrix{
X \ar[d] \ar[r] & T \ar[d] \ar[r] & \mathbf{A}^{n + m}_S \ar[ld] \\
Y \ar[r] & \mathbf{A}^n_S
}
$$
of affine schemes. By Lemma \ref{lemma-relative-pseudo-coherence-characterize}
we see that $E$ is $m$-pseudo-coherent relative to $S$ if and only if its
pushforward to $\mathbf{A}^{n + m}_S$ is $m$-pseudo-coherent. 
By Derived Categories of Schemes, Lemma
\ref{perfect-lemma-closed-push-pseudo-coherent}
we see that this is true if and only if its pushforward to $T$ is
$m$-pseudo-coherent. The same lemma shows that this holds if and only
if the pushforward to $\mathbf{A}^n_S$ is $m$-pseudo-coherent.
Again by
Lemma \ref{lemma-relative-pseudo-coherence-characterize}
this holds if and only if $Ri_*E$ is $m$-pseudo-coherent relative to $S$.
\end{proof}

\end{document}
